Bird occurrence and occupancy changes across environments after hurricanes inform monitoring priorities. We analyze \nChecklists{} eBird checklists for the Puerto Rican Emerald (\Rm), 2013--2025, combining elevation and vegetation-greenness stratification, dynamic occupancy modeling, and spatial and temporal holdout prediction to examine habitat-use differences and environmental coverage under alternative record-screening policies. Overall reporting rates were 7.59\%, 4.62\% and 6.21\% before Mar\'ia, between hurricanes and after Fiona. Standardized predictions retained high reporting levels and larger absolute rebound point estimates in high-elevation, high-greenness environments. Some lowland predictions remained below their initial levels, although all stratum-specific final-versus-initial contrast intervals spanned zero. In the corrected dynamic model, mean occupancy persistence probability was 96.8\% in high-elevation, high-greenness grids and 24.2\% and 71.3\% in low- and high-greenness lowland grids. This is the probability that a previously occupied grid remains occupied in the next primary period. Greenness was more clearly associated with occupancy persistence than with colonization, indicating that low aggregate lowland reporting can coexist with higher persistence in greener grids. Prediction gains concentrated in identification of sparse low- and middle-elevation reports. Under equal checklist-screening budgets, elevation-stratified ranking increased lowland positive-report coverage from 34.6\% to 45.6\%, while reducing overall coverage from 46.7\% to 35.2\%. Prioritizing only high reporting scores therefore does not provide balanced environmental coverage. Monitoring should retain fixed upland revisits, cover contrasting lowland greenness conditions, and use within-stratum ranking to identify supplementary candidates. Later repeat reports should be recorded separately from occupancy persistence and colonization estimated under imperfect detection, allowing environmental differences and their temporal changes to be assessed.
